\textbf{1. Kiến thức cần nhớ}
\begin{itemize}
\item \textbf{Tứ giác} $ là hình gồm bốn đoạn thẳng , BC, CD, DA$, trong đó bất kì hai đoạn thẳng nào cũng không cùng nằm trên một đường thẳng.
\item \textbf{Tứ giác lồi} là tứ giác luôn nằm trong một nửa mặt phẳng có bờ là đường thẳng chứa bất kì cạnh nào của tứ giác. (Trong chương trình, khi nói đến tứ giác mà không giải thích gì thêm, ta hiểu đó là tứ giác lồi).
\item \textbf{Định lí tổng các góc của một tứ giác}: Tổng các góc của một tứ giác bằng ^\circ$:
80669\boxed{\widehat{A} + \widehat{B} + \widehat{C} + \widehat{D} = 360^\circ}80669
\item \textbf{Góc ngoài của tứ giác}: Góc ngoài của tứ giác là góc kề bù với một góc của tứ giác. Tổng các góc ngoài tại $ đỉnh của một tứ giác lồi luôn bằng ^\circ$.
\end{itemize}

\textbf{2. Phương pháp giải các dạng}

%% DANG 8H3?10-1
\textbf{Dạng 1. Nhận biết tứ giác, các yếu tố của tứ giác}
\begin{enumerate}
\item \textbf{Xác định đỉnh, cạnh, đường chéo}:
\begin{itemize}
\item Hai đỉnh kề nhau: cùng thuộc một cạnh (ví dụ $ và $). Hai đỉnh đối nhau: không cùng thuộc cạnh nào ($ và $, $ và $).
\item Đoạn thẳng nối hai đỉnh đối nhau là \textit{đường chéo} (, BD$).
\item Hai cạnh kề nhau: cùng xuất phát từ một đỉnh. Hai cạnh đối nhau: không có điểm chung.
\end{itemize}
\item \textbf{Nhận biết tứ giác lồi}: Lấy thước kẻ đặt trùng với từng cạnh, nếu toàn bộ hình chỉ nằm về một phía của mép thước thì đó là tứ giác lồi.
\end{enumerate}

%% DANG 8H3?10-2
\textbf{Dạng 2. Tính góc dựa vào tổng các góc của tứ giác}
\begin{enumerate}
\item Sử dụng định lí: $\widehat{A} + \widehat{B} + \widehat{C} + \widehat{D} = 360^\circ$.
\item Khi biết $ góc, muốn tính góc thứ tư:
80669\widehat{D} = 360^\circ - (\widehat{A} + \widehat{B} + \widehat{C}).80669
\item Khi các góc cho dưới dạng tỉ lệ hoặc chứa ẩn $ (ví dụ $\widehat{A}:\widehat{B}:\widehat{C}:\widehat{D} = 1:2:3:4$):
\begin{itemize}
\item Biểu diễn các góc theo một biến $ (ví dụ $\widehat{A}=x, \widehat{B}=2x, \widehat{C}=3x, \widehat{D}=4x$).
\item Lập phương trình tổng bằng ^\circ$, giải tìm $ rồi suy ra số đo từng góc.
\end{itemize}
\item Sử dụng tính chất góc kề bù để tính góc ngoài hoặc góc trong tương ứng: $\widehat{A} + \widehat{A}_{\text{ngoài}} = 180^\circ$.
\end{enumerate}
